\documentclass[11pt]{article}
\usepackage[T1]{fontenc}
\usepackage{lmodern}
\usepackage[margin=26mm]{geometry}
\usepackage{amsmath,amssymb,amsthm,mathtools}
\usepackage{microtype,booktabs,longtable,array,xcolor}
\usepackage{enumitem,fancyhdr,listings}
\usepackage[colorlinks=true,linkcolor=blue!45!black,citecolor=blue!45!black,urlcolor=blue!45!black]{hyperref}
\usepackage{xurl}
\hypersetup{pdftitle={A substitution-stable seven-rule presentation of intuitionistic primitive recursive arithmetic},pdfauthor={Research prepared for Vladimir Reshetnikov},pdfsubject={Diagonal guard, exact native-kernel equivalence, and whole-domain bi-interpretation with iPRA},pdfkeywords={iPRA, primitive recursive arithmetic, intuitionistic logic, proof checker, bi-interpretation}}
\setlength{\headheight}{14pt}
\pagestyle{fancy}
\fancyhf{}
\fancyhead[L]{\small Concise iPRA: diagonal-guard introduction}
\fancyhead[R]{\small Research construction}
\fancyfoot[C]{\thepage}
\newtheorem{theorem}{Theorem}[section]
\newtheorem{lemma}[theorem]{Lemma}
\newtheorem{proposition}[theorem]{Proposition}
\newtheorem{corollary}[theorem]{Corollary}
\theoremstyle{definition}
\newtheorem{definition}[theorem]{Definition}
\theoremstyle{remark}
\newtheorem{remark}[theorem]{Remark}
\newcommand{\ipra}{\mathsf{iPRA}}
\newcommand{\FV}{\operatorname{FV}}
\newcommand{\one}{S0}
\lstset{basicstyle=\ttfamily\scriptsize,columns=fullflexible,breaklines=true,keepspaces=true,showstringspaces=false,frame=single,rulecolor=\color{black!20},xleftmargin=2mm,xrightmargin=2mm}
\setlist{topsep=4pt,itemsep=2pt}
\title{A substitution-stable seven-rule\texorpdfstring{\\}{ }presentation of intuitionistic\texorpdfstring{\\}{ }primitive recursive arithmetic}
\author{A research construction for Vladimir Reshetnikov}
\date{November 11, 2026}
\begin{document}
\maketitle
\begin{abstract}
The native scalar kernel in the current ConciseIPRA project has eight
primitive inference schemes. We replace its two introduction schemes by
one, retaining the complete original term and formula language and the
ordinary meanings of both binders. An atomic matrix receives its
progress hypothesis exactly when it has the syntactic form
\(I(y,y;u.h)=I(y,y;u.k)\); every other matrix receives its zero
instance. The zero instance of every recognized matrix is provable by
computation. This recovers all case introductions. A diagonal lift of an
arbitrary equation then recovers all equation induction. The resulting
seven-rule calculus derives exactly the original sequents and consequently
inherits the whole-domain bi-interpretation with full first-order iPRA.
The guard test is both preserved and reflected by capture-free substitution,
so the single rule supports the usual direct structural induction on proofs.

The matched, fully typed standalone Python checker has 3,244 normalized
bytes, 1,069 significant tokens and 1,124 AST nodes, compared with
3,382, 1,101 and 1,160 for the eight-rule control. This is a qualified
improvement: the existing witness-selected seven-rule control is smaller
at 3,205 bytes but its primitive mode is not stable under substitution.
The construction does not beat the 521-byte compact profiles listed in the pinned README
and does not establish absolute minimality. The new mathematics is proved
in writing; the accompanying Python certificates are checked, but no new
Lean or Rocq verification is claimed.
\end{abstract}
\tableofcontents
\newpage

\section{Result, comparison, and scope}
\label{sec:result-scope}

The comparison is frozen at ProveIt commit
\texttt{6bf83d8765614b0af51b35747e07fb057d6b2f42}, in
\texttt{Logic/PrimitiveRecursiveArithmetic/ConciseIPRA}~\cite{repo}.
Its README lists thirty-eight research packages, and its Lean development
proves the equivalence of the native scalar calculus, its ordinary
first-order realization, and iPRA~\cite{native,bridge}.

The result here is a new combination of properties: seven primitive rules,
the unrestricted original formula grammar, the unchanged conditional
readings of the binders, and an introduction rule whose mode is invariant
under capture-free numerical substitution. These are properties of the
axiomatization, not just of a Python encoding.

\begin{center}
\small
\begin{tabular}{@{}>{\raggedright\arraybackslash}p{34mm}c>{\raggedright\arraybackslash}p{34mm}>{\raggedright\arraybackslash}p{35mm}@{}}
\toprule
Presentation & Rules & Binder meanings & Introduction mode under substitution\\
\midrule
Native kernel & 8 & Original & Preserved; two rules\\
Witness-selected kernel & 7 & Original & Can change when a witness becomes literal zero\\
Uniform-progress kernel & 7 & Changed at equations & Preserved\\
Diagonal-guard kernel & 7 & Original & Preserved and reflected\\
\bottomrule
\end{tabular}
\end{center}

The prior witness rule selects progress from the matrix root and whether
the witness is literally zero~\cite{witness}. The prior uniform-progress
construction instead changes the reading of an atomic closure and modifies
closure elimination~\cite{uniform}. The prior guarded-core construction
restricts public atomic matrices beneath the paired binder to a guarded
shape~\cite{guarded}. The present construction requires no such public
restriction, changes no ordinary binder reading, and leaves closure
elimination untouched.

\paragraph{What ``simpler'' means in this report.}
There is a strict reduction from the repository's eight-rule formalized
kernel in rule count and in all three matched source-size coordinates.
There is also a structural advantage over its witness-selected seven-rule
alternative: substitution admissibility follows by a direct proof-tree
induction without changing introduction mode. There is no claim that this
candidate dominates every existing implementation, or that seven is a
lower bound. In particular, the 521-byte compact profiles listed in the
pinned README remain unbeaten. Those profiles assume a canonical guarded public goal and admit
administrative syntax, whereas the checkers measured here validate every
term, formula, scope, context and proof node in the unrestricted native
language. These are different input contracts.

\paragraph{Target theory.}
Throughout, iPRA means full first-order intuitionistic arithmetic with a
symbol for every fixed finite primitive-recursive definition, its defining
equations, the successor separation and injectivity laws, and
quantifier-free induction, with parameters and retained assumptions.
The atomic-induction formulation in the usual PR language is equivalent
through numerical characteristics of quantifier-free formulas. This is
not merely the quantifier-free equational fragment. The conventional
first-order usage is also discussed by Salehi~\cite{salehi}.

@@CORE@@

\section{Ordinary logic and the bi-interpretation with iPRA}
\label{sec:interpretation}

The equality of complete consequence relations established above is what
allows the existing arithmetic bridge to be reused. Agreement on closed
numerals or on equational theorems alone would not suffice.

\subsection{Ordinary realization and explicit logical content}

Let \(T_I\) be intuitionistic first-order logic with actual equality, the
scalar terms of the new calculus, its two iterator computation schemes,
and the following additional arithmetic principles:
\[
 \forall x\,(x=0\lor\exists z\,x=Sz),\qquad
 (0=S0)\to\bot,
\]
together with contextual induction for literal equations. These are
principles of the ordinary realization; in the native seven-rule calculus
they are derived, as are the usual logical inference rules.
The realization is
\[
 (s=t)^\circ=(s=t),\quad
 (Q_y B)^\circ=B^\circ(0)\to\forall y\,B^\circ(y),
\]
\[
 (D_{x,y}B)^\circ=
 \exists x\bigl(B^\circ(x,0)\to\forall y\,B^\circ(x,y)\bigr).
\]
All interpretation domains are the whole numerical domain, defined by
\(x=x\). No quotient equality is used.

For clarity, there are no unstated primitive logical rules in the new
calculus. Their native definitions may be written using the selector
\(C(v,a,b)=I(v,a;u.b)\), a formula selector \(M_v(A,B)\), and a formula
gate \(G_vA\). At equations the selector applies \(C\) to corresponding
sides and the gate has zero on both sides at its zero branch. At binders
these operations recurse through the formula; unmatched binder depths are
padded using an equivalent dummy-witness gated formula. First replace any
primitive \(Q_yB\) by the interderivable \(D_{x,y}B\) with an unused
witness coordinate \(x\); the case-derived closure rules prove this in
both directions. The explicit recursive definitions on the resulting
Eq/D syntax and their endpoint proofs are reproduced in
Appendix~\ref{sec:choice-introduction}.

With \(p(v)=I(v,0;u.C(u,1,0))\), one may use the following full set of
macros (fresh bound variables are always chosen):
\begin{align*}
 \bot_*&=(0=S0),\\
 A\Rightarrow_* B&=D_{u,z}M_z(A,B),\\
 A\wedge_* B&=D_{u,z}G_zM_{p(z)}(A,B),\\
 \exists_*x\,A(x)&=D_{x,y}G_yA(x),\\
 A\vee_*B&=\exists_*x\,M_x(A,B),\\
 \forall_*x\,A(x)&=A(0)\wedge_*D_{u,z}G_zA(z).
\end{align*}
The variable \(u\) is unused. The case rules give all ordinary
intuitionistic introduction and elimination rules for these macros.
The progress rule supplies exactly the additional equation induction.
No excluded-middle rule for arbitrary formulas is added.

The required separation sentence in \(T_I\) must not be omitted:
positive equations alone admit a singleton model. In the native reading,
\(0=S0\) is the designated contradiction and explosion from it is
derivable by the selector equations. The ordinary separation sentence
is precisely what connects that designated equation with ordinary
first-order falsity.

\subsection{The two term interpretations}

An iterator in iPRA is defined by primitive recursion with an update that
ignores the recursion counter:
\[
 i_h(0,a,\vec p)=a,\qquad
 i_h(Sn,a,\vec p)=h(i_h(n,a,\vec p),\vec p).
\]
This is an external translation on each fixed finite term description;
there is no internal universal evaluator for all primitive-recursive
programs.

In the other direction the scalar theory first constructs predecessor,
then a faithful equality characteristic, then quantifier-free induction,
and then pairing and projections. The dependency order matters: pairing
and equality reflection cannot be used as assumptions in the predecessor
construction. Complete inherited arguments are reproduced in
Appendices~\ref{sec:hybrid-bootstrap}, \ref{sec:eic-completion} and
\ref{sec:iterator-coding}.

If \(\pi,L,H\) are the resulting pairing and projections, a fixed PR
definition
\[
 f(0,\vec p)=g(\vec p),\qquad
 f(Sn,\vec p)=h(n,f(n,\vec p),\vec p)
\]
is compiled as
\begin{align*}
 V(n,\vec p)&=I\bigl(n,\pi(0,g^*(\vec p));
 u.\pi(SL(u),h^*(L(u),H(u),\vec p))\bigr),\\
 f^*(n,\vec p)&=H(V(n,\vec p)).
\end{align*}
Equation induction proves \(L(V(n,\vec p))=n\), and the projection
identities prove the defining equations of \(f^*\).

\subsection{Both composites and consistency}

For each fixed iterator and each fixed PR definition, equation induction
proves equality of two fixed term realizations satisfying the same
defining equations.
External induction on finite descriptions therefore identifies every
operation with its double translation. The numerical composite maps are
the identity and are internally represented by \(x=y\).

For formulas, the original native-to-ordinary and ordinary-to-native
translations already have both provable round trips. The substantive
paired-binder round trip is
\[
 D_{x,y}B\quad\dashv\vdash\quad
 \exists_*x\bigl(B(x,0)\Rightarrow_*\forall_*y\,B(x,y)\bigr).
\]
The closure round trip is obtained using a dummy witness. Because the new
calculus has exactly the original sequents, every such round-trip proof
is also a proof in the new calculus. Thus the composition with the
\(T_I\)--iPRA bridge is a genuine bi-interpretation on the whole
numerical domain, not merely a mutual consistency reduction.

Effective transformations of finite derivations carry a proof of native
\(0=S0\) to an iPRA proof of \(\bot\), and conversely. This yields
equiconsistency. It asserts neither that iPRA proves its own consistency
nor that the new system has a special standard-domain semantics.

\section{Standalone checkers and reproducible measurements}
\label{sec:implementation}

The three compared checkers share the same explicit de Bruijn syntax,
ordinary JSON input, scope checks, equality rule and input/output contract.
They differ only in their introduction rules. The eight-rule control is
the repository's \texttt{one\_context/direct/native\_checker.py} at the
pinned commit; the witness and diagonal variants are matched edits of that
control. The filename of the control records its package location, not a
change to the two-template equality rule used in this comparison.

The data object has exactly four fields:
\begin{lstlisting}
{"variables": n, "context": [...], "goal": formula, "proof": proof_tree}
\end{lstlisting}
Here \(n\) is a nonnegative integer. Context entries and the goal must be
well-scoped formulas. Terms are tagged as \texttt{v}, \texttt{z},
\texttt{s}, and \texttt{i}; formulas as \texttt{=}, \texttt{q} and
\texttt{d}. Only an iterator's body binds its state variable.
The matrix of \(D\) binds \(y\) at index 0 and \(x\) at index 1.
All proof nodes are checked and synthesize their conclusions.

The new introduction is implemented by opening the witness and retaining
the opened equation itself exactly when both iterator prefixes are
\texttt{['i', ['v', 0], ['v', 0]]}; otherwise it retains the shifted
zero instance. Well-formedness is checked before testing these prefixes.
Witness substitution cannot manufacture the bound \(y\) in one of these
positions, so checking the guard after opening the witness is equivalent
to checking the original matrix.

Acceptance means exit code zero and exactly \texttt{ACCEPT} followed by a
newline. The programs use neither assertions for logical checks nor
\texttt{eval}, a proof database, an external prover, a trusted compiler,
or an imported axiom file. The untrusted proof compiler and test producer
are separate from the checker. The trusted runtime includes CPython and
its standard \texttt{json} and \texttt{sys} modules, just as the controls
do. Resource exhaustion may reject a large valid finite certificate;
no claim of fixed-machine resource completeness is made.

\paragraph{Rule correspondence.}
The implementation uses one branch for each displayed primitive rule.
By induction on a successful call to \texttt{infer}, its result is a
well-scoped formula derivable from the supplied context: the hypothesis
branch checks membership, the computation branches check their term
arguments, and transport and QE compare the exact substituted minor
premise. DE checks its continuation against a shifted conclusion under a
new variable, excluding eigenvariable escape. DI checks exactly the
single local hypothesis selected by the guard. The binder increments in
\texttt{walk} are zero for ordinary term arguments, one for an iterator
body and a Q matrix, and two for a D matrix; structural induction gives
the capture-free meanings of \texttt{shift} and \texttt{sub}.
Conversely every finite derivation of the displayed calculus has the
corresponding tagged certificate. This is a written correspondence
argument for the idealized unbounded-resource program, not a
proof-assistant-checked theorem about the Python implementation.

@@METRICS@@

The normalization applies \texttt{ast.unparse(ast.parse(source))} and
appends one line-feed byte, under CPython 3.12.14. Significant tokens exclude encoding, indentation,
dedentation, newlines, comments and the end marker. AST nodes are counted
with \texttt{ast.walk}. Informational gzip sizes use compression level 9
and \texttt{mtime=0}. The same procedure is used for all three modules;
none of these figures omits a checker helper.

The diagonal rule saves 138 normalized bytes (4.08\%), 32 tokens
(2.91\%) and 36 AST nodes (3.10\%) against the eight-rule control.
It costs 39 bytes against the witness-selected control. This is a
Pareto tradeoff between source size and a substitution-invariant primitive
interface, not a new record for unrestricted minimum source size.

\section{Verification and limitations}
\label{sec:validation}

@@VALIDATION@@

\paragraph{Executable translations.}
The untrusted translator implements both complete proof-tree translations,
including elimination of arbitrary native \texttt{qi} nodes. A separate
suite checks 29 certificate cases, with 93 checker validations covering
every source rule in both calculi and both round trips. In particular it
compiles an identity induction, an induction whose base is the unproved
equation \(a=b\), and an induction using both the current induction
hypothesis and a retained parameter equation. The three translated
proofs are also accepted through the public command-line interface in
ordinary and optimized Python modes. The count separates 87 direct calls
to the standalone checker modules from six command-line processes.

\begin{center}\small
\begin{tabular}{@{}lrr@{}}
\toprule
Certificate & Original JSON bytes & Translated JSON bytes\\
\midrule
Identity induction & 485 & 11,506\\
Conditional seed equality & 885 & 14,218\\
Contextual update congruence & 1,965 & 20,283\\
\bottomrule
\end{tabular}
\end{center}
These are certificate sizes, not checker sizes. The larger proof trees
make explicit the logical work removed from the primitive kernel. The
report claims an effective translation and reports these examples; it
does not claim a polynomial bound on serialized tree size.

The checker tests establish that the supplied finite certificates are
processed as reported. They do not establish universal soundness,
completeness, or consistency. Those claims rest on the written
metatheorems and the inherited arithmetic bridge. The pinned repository's
Lean formalization verifies the original bridge; it does not automatically
verify a new Python checker or the new diagonal-rule comparison.
No new Lean or Rocq build was performed for this delivery.

The archive includes the new construction, checker sources, proof
translation code and examples, measurements, reproduction scripts,
provenance, and the inherited arithmetic arguments with attribution.
The report's new contribution is the diagonal dispatch theorem and its
matched implementation. It does not claim a new predecessor construction,
a new PR compilation method, an absolute minimum number of rules, or a
smaller version of the 521-byte guarded-profile checker listed in the pinned README.

\section{Further research questions}
\label{sec:research}

\begin{enumerate}
\item \textbf{Lean comparison with the existing native kernel.}
Add a deeply embedded \(\mathcal G\) with the guard as a structural
predicate. Prove guard preservation and reflection first, followed by
substitution and cut directly. Then implement the equation-only
implication construction and the proof of unrestricted progress.
The target is an exact sequent theorem, not only consistency.
\item \textbf{The smallest stable guard.}
Can a cheaper decidable structural class retain all three required
properties: provable zero instances, a lift of arbitrary equations,
and reflection under substitution? The diagonal guard is a concrete
solution, not a lower-bound theorem.
\item \textbf{Proof size and sharing.}
The translations are effective, but syntactic lifting and repeated local
derivations may expand tree certificates. Determine explicit size bounds
and compare trees with a checked DAG format. Sharing must not enlarge the
trusted logical rule set implicitly.
\item \textbf{Combining independent equality reductions.}
The repository's crossed one-template transport is compatible at the
mathematical level with the present introduction theorem. A matched
implementation could determine whether their combination improves the
smallest fully typed checker while retaining a transparent proof format.
Such a measurement is separate from the results reported here.
\item \textbf{A fair compact-profile comparison.}
Can the new substitution-stable rule be encoded within a compact
administrative representation with a complete public-goal validator,
and improve its normalized source size? The current typed experiment
does not answer that question.
\end{enumerate}

\appendix
\section{Attribution and inherited proof material}
\label{sec:attribution}

The following five sections reproduce the mathematical bridge from the
pinned repository's \texttt{Research/certified\_zero/reference/} material
with light introductory wording changes. The Choice-DI theorem and its
logical macros, the ordinary realization, scalar predecessor bootstrap,
equality reflection, quantifier-free induction, pairing and PR compilation
are inherited results, not new results of this report. Exact fetched
sources and their provenance are included separately in the archive.
The repository distributes these materials under the MIT No Attribution
license; its license text is included.

@@INHERITED@@

\section{The complete new checker}
\label{sec:listing}

This is the full standalone program. Its source-size coordinates include
all the code below. The shorter normalized rendering differs only in
formatting.
\begin{lstlisting}[language=Python]
@@CHECKER@@
\end{lstlisting}

\begin{thebibliography}{10}
\bibitem{repo} ProveIt contributors, \emph{Concise presentations of
intuitionistic primitive recursive arithmetic}, README at commit
\texttt{6bf83d8765614b0af51b35747e07fb057d6b2f42}.
\url{https://github.com/VladimirReshetnikov/ProveIt/tree/6bf83d8765614b0af51b35747e07fb057d6b2f42/Logic/PrimitiveRecursiveArithmetic/ConciseIPRA}.
\bibitem{native} ProveIt contributors, \emph{Native.lean} and
\emph{NativeStruct.lean}, same commit, \texttt{Lean/ConciseIPRA/}.
\bibitem{bridge} ProveIt contributors, \emph{NativeBiInterpretation.lean}
and \emph{BiInterpretation.lean}, same commit, \texttt{Lean/ConciseIPRA/}.
\bibitem{witness} ProveIt contributors, \emph{witness\_rule} research
package and \emph{NativeOneIntro.lean}, same commit.
\bibitem{guarded} ProveIt contributors, \emph{guarded\_core} research
package, same commit. See also \emph{guarded\_source} and
\emph{deterministic\_auxiliary}.
\bibitem{uniform} ProveIt contributors, \emph{uniform\_progress} research
package, same commit.
\bibitem{erasure} ProveIt contributors, inherited
\emph{scalar\_predecessor.tex},
\path{Research/certified_zero/reference/}, same commit.
\bibitem{forward} ProveIt contributors, inherited
\emph{equality\_reflection.tex},
\path{Research/certified_zero/reference/}, same commit.
\bibitem{progress} ProveIt contributors, inherited
\emph{primitive\_recursion.tex},
\path{Research/certified_zero/reference/}, same commit.
\bibitem{salehi} S. Salehi, \emph{Decidable formulas of intuitionistic
primitive recursive arithmetic}, Reports on Mathematical Logic 36
(2002), 55--61.
\url{https://rml.tcs.uj.edu.pl/rml-36/36-SALEHI.pdf}.
\end{thebibliography}
\end{document}
